\subsection{A nervous-system test: content support versus total arrival time}\label{sec:nervous-test}
A particularly relevant public candidate is the release accompanying \citet{michaels2025}. Its Dryad description lists session-level spike timing, hand and joint kinematics, and trial timing and conditions, with a recording-session information table \citep{michaelsdata2025}. The experiment uses mechanically perturbed sensorimotor behavior and neural population recordings. It can support an initial held-out test of whether matched perturbation contrasts occupy different receiver readouts and arrive at different times. Mechanical perturbation is a timed sensory input, not a selective intervention on an isolated cortical edge.

\paragraph{What the release actually covers.} A fresh audit of the official Dryad version manifest identifies 49 session files (29 from monkey M and 20 from monkey P). The paper's Extended Data Table 1 provides a more restrictive eligibility screen: among S1, M1, PMd, and dlPFC, there are nine sessions with both PMd and dlPFC, three with S1 and M1, two with M1 and PMd, and two with S1 and PMd. These are simultaneous-area candidates, not 49 independent source--receiver pairs. No listed session pairs dlPFC directly with M1. Several PMd--dlPFC sessions contain very few PMd neurons, so a frozen minimum-feature/quality screen may reduce the nine eligible sessions further. The transcribed area-pair table and official file manifest accompany the code. Raw trial synchrony and preprocessing still require file-level verification; metadata-file download endpoints returned access errors in this audit, and no new waveform fit was performed.

The proposed analysis first freezes the source and receiver features, readouts, preprocessing delays, and training/test split. Next it compares three models on identical held-out trials: (a) a common scalar delay with fixed projection; (b) content-specific projected routes with their independently constrained delays; (c) a flexible delayed nonlinear predictor using the same inputs. A success for (b) must predict both a selective near-null contrast and the later supported contrast, improve paired held-out error with session-level uncertainty, and survive the flexible rival. Inferring routes additionally requires simultaneous source--target coverage and controlled alternatives. Session metadata must establish that coverage before any paired-population model is fitted. Conditioning on movement or trial outcome can itself bias interpretation.

A decisive perturbation of the proposed detour would change its later response while preserving an independently supported early route. A paired-site peripheral stimulation study is a complementary timing calibration: subtracting proximal and distal onset times can cancel shared downstream stages under a justified common-path model. It does not identify a content translation from muscle voltage alone. The open PhysioNet EMG examples contain voluntary muscle recordings, rather than synchronized paired-site conduction stimulation and source/receiver population contrasts \citep{emgexamples2009}; they are therefore unsuitable for testing the complete model. F-TRACT's downloadable derived atlas adds a model-dependent latency reference \citep{ftractzenodo2022}, not independently calibrated content maps. None of these resources alone identifies the entire nervous-system network, its receptive phase, its interaction rank, and its learning rule.

No additional biological waveform was fitted in this revision. The prospective contribution is the jointly calibrated prediction in Theorem~\ref{thm:network}, not retrospective agreement obtained by interpreting any observed latency as its delay term. Sample planning must use independent-session variance and the number of tested source contrasts; repeating dense within-trial time points does not increase biological replication.

\paragraph{Concrete planning and rejection.} Let $D_s$ be the held-out MSE of a rival minus that of the projected-route model in independent replication unit $s$, and $d=\mathbb E D_s/\operatorname{SD}(D_s)$. For two prespecified rival comparisons, use two-sided level $0.025$ each (Bonferroni familywise level $0.05$). Under normal independent unit-level differences, exact noncentral-$t$ planning gives:
\begin{center}\begin{tabular}{@{}ccc@{}}\toprule
Assumed standardized advantage $d$ & Units for 80\% power & Units for 90\% power\\\midrule
0.3 & 109 & 141\\
0.5 & 41 & 53\\
0.8 & 18 & 23\\\bottomrule
\end{tabular}\end{center}
These are assumed-effect calculations, not estimates from the release. Achieving 90\% power for each comparison guarantees at least 80\% joint power by the union bound. Freeze a practically meaningful MSE margin as well as the comparisons; a tiny significant improvement is insufficient. Repeated sessions from the same animal require clustering, and neurons and time bins are not independent replication units. Two animals cannot establish population-level uncertainty across animals from session count alone.

A selective near-null prediction needs an equivalence test rather than a nonsignificant response. As a concrete planning illustration, two one-sided tests at level $0.025$, a margin of $0.3$ unit-level response SD, true mean zero, and independent normal contrasts require 119 units for 80\% equivalence power or 147 for 90\%. These values integrate over the estimated sample variance, rather than treating it as known. Prespecify an early window, a null contrast, a supported contrast, and a minimum late-response amplitude from independent calibration. Joint confirmation requires the early contrast within its equivalence margin and the later contrast above its positive-effect bound; correct for any additional contrasts or windows. These planning assumptions may be unattainable in this release, which would make it exploratory for the joint claim.

With adequate calibration and precision, an early-response confidence interval excluding the predicted near-null range, a late-response upper confidence bound below the calibrated minimum, or timing inconsistent with the independently constrained route excludes the fitted route hypothesis. Failure to outperform a flexible rival limits predictive usefulness; it does not by itself disprove the algebraic network theorem. Without resolved source--receiver coverage or enough independent replication, the result is inconclusive. Supplement S20 and \texttt{verify\_revision\_planning.py} reproduce these calculations.
